\section*{Bab \ref{ch:g9:chromosomes} --- Kromosom}

\begin{solution}{exo:g9:chromosomes:1}
Sebuah benda mirip benang di dalam \omterm{def:g6:cell-unit-of-life:parts}{inti selnya} --- isinya
yang menggulung untuk diangkut. Terlihat, bila diwarnai
baik, ketika \omterm{def:g5:first-look-at-cells:cell}{selnya} membelah; terurai dan tak terlihat di
antara pembelahannya.
\end{solution}

\begin{solution}{exo:g9:chromosomes:2}
46 pada setiap \omterm{def:g5:first-look-at-cells:cell}{sel} tubuh; terpilah sebagai 23 pasang; 23 ---
satu mitra tiap pasang --- di dalam \omterm{def:g8:sexual-reproduction:gametes}{gametnya}.
\end{solution}

\begin{solution}{exo:g9:chromosomes:3}
$23 + 23 = 46$: \omterm{def:g5:flower-to-fruit:steps}{pembuahannya} memulihkan pasangannya, satu
mitra tiap pasang dari ibunya, satu dari ayahnya --- kedua
orang tuanya ada di setiap \omterm{def:g5:first-look-at-cells:cell}{sel} tubuh anaknya.
\end{solution}

\begin{solution}{exo:g9:chromosomes:4}
\omterm{def:g8:sexual-reproduction:gametes}{Gamet} yang tidak dibagi dua akan melipatduakan tumpukannya
tiap angkatan: gamet-46 membuat anak-92, lalu cucu-184.
Pembagian duanya menjaga \omterm{def:g5:biodiversity:species}{spesiesnya} tetap 46.
\end{solution}

\begin{solution}{exo:g9:chromosomes:5}
XX (\omterm{def:g9:chromosomes:chromosomes}{kariotipe} seorang anak perempuan) dan XY (\omterm{def:g9:chromosomes:chromosomes}{kariotipe}
seorang anak laki-laki). \omterm{def:g8:sexual-reproduction:gametes}{Gamet} ayahnya: \omterm{def:g2:animals-grow-up:born}{telurnya} semua
mengangkut X, spermanya mengangkut X atau Y separuh-separuh
--- sperma yang tibalah yang memutuskan.
\end{solution}

\begin{solution}{exo:g9:chromosomes:6}
Perangkat \omterm{def:g9:heredity-and-traits:traits}{sifat} yang dikenali sebagai sindrom Down. Ia
membuktikan bahwa benangnya mengangkut petunjuk \omterm{def:g9:heredity-and-traits:traits}{sifatnya}:
bilangan satu benang berubah, banyak \omterm{def:g9:heredity-and-traits:traits}{sifat} bergeser.
\end{solution}

\begin{solution}{exo:g9:chromosomes:7}
Pada pembelahan pembagi duanya: mitra yang mana dari tiap
pasang yang diterima sebuah \omterm{def:g8:sexual-reproduction:gametes}{gamet} diputuskan secara bebas,
pasang demi pasang --- $2^{23}$ setengah-tumpukan yang
mungkin tiap orang tua, sebelum \omterm{def:g5:flower-to-fruit:steps}{pembuahannya} mengalikan dua
ruang semacam itu.
\end{solution}

\begin{solution}{exo:g9:chromosomes:8}
Penentunya, bukan \omterm{def:g9:heredity-and-traits:traits}{sifatnya}: benangnya mengangkut petunjuk
--- ringkas, terbawa \omterm{def:g5:first-look-at-cells:cell}{sel}. Diangkut tanpa diperlihatkan:
kedua mitranya menumpang di setiap \omterm{def:g5:first-look-at-cells:cell}{sel} apa pun yang
ditampilkan. Pilihan yang berbeda dari anak ke anak:
pembagian duanya membagikan satu mitra tiap pasang, secara
bebas pada tiap kesempatan.
\end{solution}

\begin{solution}{exo:g9:chromosomes:9}
Tiap \omterm{def:g5:flower-to-fruit:steps}{pembuahan} adalah lemparan yang baru: \omterm{def:g8:reproduction-and-environment:population}{populasi} spermanya
separuh X, separuh Y pada tiap percobaan, dan tidak ada
\omterm{def:g2:born-grow-die:cycle}{kelahiran} lampau yang menjangkau perlombaannya. Tiga anak
perempuan adalah tiga lemparan yang adil --- yang keempat
sama adilnya.
\end{solution}

\begin{solution}{exo:g9:chromosomes:10}
(a) \omterm{def:g5:first-look-at-cells:cell}{Sel} tubuh seorang anak perempuan atau perempuan dewasa:
tumpukan penuh, XX. (b) Sebuah \omterm{def:g8:reproductive-systems:male}{sel sperma} yang mengangkut Y:
sebuah \omterm{def:g8:sexual-reproduction:gametes}{gamet}, setengah tumpukan --- \omterm{def:g5:flower-to-fruit:steps}{pembuahannya} akan
menghasilkan anak laki-laki. (c) \omterm{def:g5:first-look-at-cells:cell}{Sel} tubuh dengan kelebihan
pada pasang 21: 47 benang --- \omterm{def:g9:chromosomes:chromosomes}{kariotipe} sindrom Down, XX:
milik seorang anak perempuan.
\end{solution}

\begin{solution}{exo:g9:chromosomes:11}
Setiap \omterm{def:g2:animals-grow-up:born}{telur} yang dipasok para ratu mengangkut sebuah X ---
sumbangan mereka pada pasang 23 tidak mengizinkan pilihan
lain. Keragaman X-atau-Y menumpang di dalam \omterm{def:g8:sexual-reproduction:gametes}{gamet} para raja
saja: bila catatannya harus menetapkan tempat koin
lemparannya, ia duduk di dalam sperma rajanya --- dan
bagaimanapun lemparannya adalah peluang, yang tak dapat
disalahkan kepada siapa pun.
\end{solution}

\begin{solution}{exo:g9:chromosomes:12}
Teka-tekinya: satu tumpukan, banyak jenis \omterm{def:g5:first-look-at-cells:cell}{sel} --- bila
setiap \omterm{def:g5:first-look-at-cells:cell}{sel} memuat 46 benang petunjuk yang sama, mengapa
sebuah \omterm{def:g8:nervous-communication:neuron}{neuron} bukan sebuah \omterm{def:g5:first-look-at-cells:cell}{sel} \omterm{def:g1:the-five-senses:sense}{kulit}? Arah pemecahannya: \omterm{def:g5:first-look-at-cells:cell}{sel}
pastilah berbeda dalam petunjuk mana yang \emph{dibacanya},
bukan dalam petunjuk mana yang dimuatnya --- tumpukannya
adalah sebuah perpustakaan, dan pengkhususannya adalah
sebuah kebijakan peminjaman. Apa satuan yang dapat dibaca
pada sebuah \omterm{def:g9:chromosomes:chromosomes}{kromosom} adalah pertanyaan bab berikutnya.
\end{solution}

\begin{solution}{pb:g9:chromosomes:1}
\textbf{1.} A: \omterm{def:g5:first-look-at-cells:cell}{sel} tubuh seorang anak laki-laki atau
laki-laki dewasa --- 46, XY. B: \omterm{def:g5:first-look-at-cells:cell}{sel} tubuh seorang anak
perempuan atau perempuan dewasa --- 46, XX. C: sebuah \omterm{def:g8:sexual-reproduction:gametes}{gamet}
--- 23, satu mitra tiap pasang; X-nya membuatnya entah
sebuah \omterm{def:g4:animal-reproduction:sexual}{sel telur} entah sebuah sperma-X.

\textbf{2.} Sebuah \omterm{def:g8:sexual-reproduction:gametes}{gamet}. Angka 23-nya mencatat pembelahan
pembagi dua yang khusus: dari sebuah sel-46, satu mitra tiap
pasang diambil, tidak ada pasang yang lengkap.

\textbf{3.} 47 benang, kelebihannya pada pasang 21, pasang
ke-23 XX: \omterm{def:g9:chromosomes:chromosomes}{kariotipe} seorang anak perempuan dengan sindrom
Down --- satu benang salah hitung, perkembangan yang bergeser.

\textbf{4.} \omterm{def:g8:sexual-reproduction:gametes}{Gamet} ayahnya --- X-atau-Y spermanya. Peluang
pada pertemuannya: satu dari dua, ke tiap arah.

\textbf{5.} Karena tubuhnya tumbuh dari satu \omterm{prop:g8:sexual-reproduction:core}{zigot} lewat
pembelahan penyalin: tiap pembelahan menggandakan
tumpukannya lalu membagikan satu salinan penuh kepada tiap
\omterm{def:g5:first-look-at-cells:cell}{sel} anaknya --- 46 salinan yang setia, dari \omterm{def:g1:the-five-senses:sense}{kulit} sampai
\omterm{def:g8:nervous-communication:neuron}{neuron}.

\textbf{6.} Pembelahan pembagi dua pada \omterm{def:g8:sexual-reproduction:gametes}{gametnya}, dengan
aturannya: satu mitra dari tiap ke-23 pasangnya --- sebuah
setengah tumpukan yang lengkap, tanpa pasang yang berganda,
tanpa yang tercecer.

\textbf{7.} Bahwa anaknya menerima sebuah X darinya: mitra
mana pun yang tiba dari sisi yang lain, sumbangan C pada
pasang 23 sudah dipatok --- gamet-X menuliskan X ke dalam
pasang ke-23 anaknya tanpa syarat.

\textbf{8.} Masing-masing dari ke-23 pasangnya menyumbangkan
salah satu mitranya, secara bebas: $2 \times 2 \times \dots$,
23 kali --- $2^{23}$ setengah tumpukan, beberapa juta, dari
satu orang tua. \omterm{def:g5:flower-to-fruit:steps}{Pembuahannya} memasangkan satu pembagian
semacam itu dari tiap orang tua: kedua ruangnya mengalikan
diri, dan hitungan pengocokannya menjelaskan mengapa
adik-kakak tidak pernah berulang.

\textbf{9.} Tidak ada yang dilakukan salah: kelebihan 21-nya
timbul pada sebuah kecelakaan pembagian dua --- kedua mitra
sepasang gagal memisahkan diri --- sebuah kemalangan
mekanismenya, bukan kemalangan perbuatan atau harga diri
orang tuanya. Benang ke-47 itu menggeser perkembangan
anaknya; ia tidak mengubah apa pun tentang anaknya yang
sepenuhnya anak mereka --- kebenaran konselingnya dan
kebenaran sitogenetikanya sepakat.

\textbf{10.} Satu dari dua, persis seperti pada setiap
\omterm{def:g2:born-grow-die:cycle}{kelahiran}: \omterm{def:g8:reproduction-and-environment:population}{populasi} spermanya separuh X, separuh Y pada tiap
\omterm{def:g5:flower-to-fruit:steps}{pembuahannya}, dan lemparan yang lalu tidak meninggalkan
jejak padanya. Deretannya terasa bermakna; mekanismenya
tidak menyimpan ingatan.

\textbf{11.} Keputusannya dibuat pada \omterm{def:g5:flower-to-fruit:steps}{pembuahannya}, oleh
sperma mana yang tiba --- di dalam sumbangan ayahnya, pada
milimeter terakhir perlombaannya. Sebuah pola makan tidak
menjangkau pembuatan separuh-separuh \omterm{def:g8:reproduction-and-environment:population}{populasi} spermanya,
tidak pula hasil perlombaannya: klaimnya tidak punya mata
rantai untuk dihadang atau dicondongkan ---
\cref{met:g8:contraception:evaluate} pertanyaan 1, gagal lagi.

\textbf{12.} \omterm{def:g9:heredity-and-traits:heredity}{Pewarisan sifatnya} menumpang pada benang yang
dapat dihitung di dalam \omterm{def:g6:cell-unit-of-life:parts}{inti selnya}: dibagi dua ke dalam
\omterm{def:g8:sexual-reproduction:gametes}{gametnya}, dipulihkan pada \omterm{def:g5:flower-to-fruit:steps}{pembuahannya}, menggeser \omterm{def:g9:heredity-and-traits:traits}{sifat} bila
salah hitung, memutuskan kelamin pada pasang 23 --- empat
foto, satu kendaraan.
\end{solution}
